\clearpage
\section{Apendice: mapa pagina por pagina de los slides}
Este apendice lista cada pagina/slide de la carpeta fuente y su contenido principal. Varias paginas son animaciones repetidas del mismo slide; se conservan para trazabilidad.
\subsection*{CS3014\_\_\_Strings.pdf (23 paginas)}
\begin{description}[leftmargin=!,labelwidth=4.2em,style=nextline]
\item[Pag. 1] CS3014 - Estructuras de Datos Avanzadas Notas de clase 01 Strings 31 de mayo de 2026 Prof. V'  ctor Racs' o Galv' an Oyola 1. Notaci' on Sea Sigma un conjunto ordenado de s'  mbolos (a partir de ahora,caracteres), el cual llamaremos alfabeto. Diremos que una secuencia finitaSque consiste de caracteres de Sigma es unacadenasobre el alfabeto Sigma. Por convenci' on, denot...
\item[Pag. 2] CS3014 - Estructuras de Datos Avanzadas Notas de clase 01 SyTdifieren en una posici' on. Seapla menor posici' on en la que difieren, entonces S[p]<T[p]. Por convenci' on, usaremos comparaci' on lexicogr' afica al comparar dos cadenas. 2. Suffix Array El suffix array es una estructura de indexing cuya caracter'  stica principal es la poca cantidad de memoria...
\item[Pag. 3] CS3014 - Estructuras de Datos Avanzadas Notas de clase 01 Algoritmo 1:Construir-Mapeos(S, n) 1k<-1 +floor(log 2 n); 2Seamapeo[k][n] un arreglo de enteros ; 3p<- \{0, . . . , n-1\}; 4OrdenarpporS i creciente ; 5Asignarmapeo[0] de acuerdo ap; 6parad<-1ak-1hacer 7v<- \{(mapeo[d-1][i], mapeo[d-1][i+ 2 d-1]) : 0<=i<=n-2 d\}; 8Filtrarppara quedarnos con solo losi in [0, n-2 d] ; 9O...
\item[Pag. 4] CS3014 - Estructuras de Datos Avanzadas Notas de clase 01 Niveld= 1 Se filtra el vectorpy mapeamosi7->i-2 d-1 =i-1: (1,4,7,10,8,9,2,3,5,6)- ->(0,3,6,9,7,8,1,2,4,5). Luego de ordenar de manera estable seg' unmapeo[0], se obtiene: p= (7,1,4,0,9,8,3,6,2,5). El nuevo mapeo es: i 0 1 2 3 4 5 6 7 8 9 10 mapeo[1] 2 1 6 5 1 6 5 0 4 3 ? Niveld= 2 Se filtrapy mapeamosi7...
\item[Pag. 5] CS3014 - Estructuras de Datos Avanzadas Notas de clase 01 haremos ... ?o s'  ? En realidad es posible extenderSusando 2 m -1 caracteres sentinelas \$. Esta ' ultima condici' on es necesaria para representar cadenas que originalmente no estar'  an completas (en cuyo caso, la relaci' on cuando una es prefijo de otra se establece correctamente ya que \$ precede a...
\item[Pag. 6] CS3014 - Estructuras de Datos Avanzadas Notas de clase 01 Iteraci' onlen= 1 El mapeo inicial queda dado por: i 0 1 2 3 4 5 6 7 8 9 10 11 mapeo 2 1 4 4 1 4 4 1 3 3 1 0 Iteraci' onlen= 2 Luego de asignar el vectorby second, este queda de la siguiente manera: i 0 1 2 3 4 5 6 7 8 9 10 11 by second 2 1 4 4 1 4 4 1 3 3 1 0 Y losheadquedar'  an de la siguiente mane...
\item[Pag. 7] CS3014 - Estructuras de Datos Avanzadas Notas de clase 01 i 0 1 2 3 4 5 6 7 8 9 10 11 head 0 1 2 3 5 6 7 8 9 10 11 0 Luego de ejecutar la l'  neas 10-13 del algoritmo, se obtiene: a= (11,10,7,4,1,0,9,8,6,3,5,2). El nuevo mapeo es: i 0 1 2 3 4 5 6 7 8 9 10 11 mapeo 5 4 11 9 3 10 8 2 7 6 1 0 Iteraci' onlen= 16 Luego de asignar el vectorby second, este queda de...
\item[Pag. 8] CS3014 - Estructuras de Datos Avanzadas Notas de clase 01 Rk = Suf(S, k) agrupado en grupos de 3 caracteres consecutivos Consideremos que si alg' un grupo contiene menos de 3 caracteres, completaremos dichos caracteres con sentinelas \$. Por ejemplo, paraS=mississippiconn= 11, se dar'  a que R1 =(iss)(iss)(ipp)(i\$\$) R2 =(ssi)(ssi)(ppi) Adem' as, se defineR=R...
\item[Pag. 9] CS3014 - Estructuras de Datos Avanzadas Notas de clase 01 i in B 0 yj in B 1: Comparamos por pares como enB 0. i in B 0 yj in B 2: Comparamos por ternas (Si, Si+1,rank S(i+2)) y (Sj, Sj+1,rank S(j+ 2)). Asumiremos que tendremos una funci' on llamadaCompare-DC3(i, j, S, rnk) que permita comparar las posicionesiyjseg' un este criterio (asumiremos que devuelve una respuest...
\item[Pag. 10] CS3014 - Estructuras de Datos Avanzadas Notas de clase 01 Algoritmo 5:DC3(n, S) 1SeanB k =\{0<=i < n:i=k(m' od 3)\}parak= 0,1,2 ; 2sin=1 (m' od 3)entonces 3B 1 <-B 1  U  \{n\}; 4C<-B 1  U B 2 ; 5OrdenarCpor (S[x], S[x+ 1], S[x+ 2]) conRadix-Sort; 6Crearsigma:C7->Z + 0 seg' un el orden deC; 7R<- 2Q i=1 Q x in Bi sigma((S[x], S[x+ 1], S[x+ 2])) ; 8A R <-DC3(|R|, R) ; 9Searnk[0. . . n]...
\item[Pag. 11] CS3014 - Estructuras de Datos Avanzadas Notas de clase 01 Algoritmo 6:Classify(n, S) 1SeaT[0. . .(n-1)] un arreglo de tipos ; 2T[n-1]<-S; 3parai<-n-2a0hacer 4siS[i] =S[i]entonces 5T[i]<-T[i+ 1] ; 6en otro caso 7siS[i]< S[i+ 1]entonces 8T[i]<-S; 9en otro caso 10T[i]<-L; 11devolverT; Debido a la definici' on de los tipos, se puede probar el siguiente lema: Lema 7 S...
\item[Pag. 12] CS3014 - Estructuras de Datos Avanzadas Notas de clase 01 2. Iteramos de izquierda a derecha (derecha a izquierda) sobre los elementos deA, siS A[i]-1 es de tipoL(S), entonces lo colocamos al frente (final) de su bloque y movemos el puntero correspondiente. Es sencillo notar que la complejidad del algoritmo anterior esO(n), ya que podemos mantener el inicio...
\item[Pag. 13] CS3014 - Estructuras de Datos Avanzadas Notas de clase 01 b d e f d f d f b d e d d a f d\$ Tipo S S S L S L S L S S L L L S L L S/L Cuadro 1: Tipos de sufijo de la cadena "bdefdfdfbdeddafd\$" Subcadenas S Subcadenas L afd\$ bd de dedda dfb dfd efd d\$ daf dd ed fbde fd fdf Cuadro 2: Subcadenas de la cadena "bdefdfdfbdeddafd\$" Si asoci' aramos cada subcadena a s...
\item[Pag. 14] CS3014 - Estructuras de Datos Avanzadas Notas de clase 01 Existe una posici' onjtal queX j  =Y j yX j <Y j. Y(X) es prefijo propio deX(Y). A partir de este enfoque, se cumple el siguiente lema. Lema 13 SeaS ' la cadena compuesta por los mapeos de las subcadenas elegidas deS, entonces para dos sufijosS i yS j con sus posiciones correspondientes enS ' i' yS '...
\item[Pag. 15] CS3014 - Estructuras de Datos Avanzadas Notas de clase 01 funciones deben modificarse. Algoritmo 9:Group-By-Character(n, S) 1SeaA[n] un arreglo de enteros ; 2Seahead[n] un arreglo de enteros inicializado en 0 ; 3parai<-0an-1hacer 4head[S[i]]<-head[S[i]] + 1 ; 5sum<-0 ; 6parai<-0an-1hacer 7sum<-sum+head[i] ; 8head[i]<-sum-head[i] ; 9parai<-0an-1hacer 10A[head[S[i]]]...
\item[Pag. 16] CS3014 - Estructuras de Datos Avanzadas Notas de clase 01 Algoritmo 12:Move-To-Front(R, lptr, L j, l, r) 1parai<-lar-1hacer 2start<-L j[i]-j; 3pos bucket<-R[start] ; 4f ront bucket<-lptr[pos bucket] ; 5sif ront bucket <0entonces 6f ront bucket<- not f ront bucket; 7lptr[pos bucket]<- not (f ront bucket+ 1) ; 8sif ront bucket=pos bucketentonces 9lptr[pos bucket]<- not lptr[p...
\item[Pag. 17] CS3014 - Estructuras de Datos Avanzadas Notas de clase 01 posiciones con sufijos de tipo S. Algoritmo 15:Build-Reduced-String-With-S(n, S, T) 1A<-Group-By-Character(n, S) ; 2Dist<-Get-S-Distances(n, S, T) ; 3m<- n-1 m' ax i=0 \{Dist[i]\}; 4SeanL 1 . . . Lm listas vac'  as de enteros ; 5parai<-0an-1hacer 6idx<-A[i]; 7siDist[idx]>0entonces 8Agregaridxal final deL Dis...
\item[Pag. 18] CS3014 - Estructuras de Datos Avanzadas Notas de clase 01 Algoritmo 16:SAKA(n, S) 1sin= 1entonces 2devolver\{0\}; 3en otro caso 4T<-Classify(n, S) ; 5cnt L <- |\{0<=i < n-1 :T i =L\}|; 6cnt S <-n-1-cnt L ; 7sicnt L < cnt S entonces 8S ' <-Build-Reduced-String-With-L(n, S, T) ; 9siS ' tiene todos sus caracteres diferentesentonces 10CalcularXen funci' on al caracter de...
\item[Pag. 19] CS3014 - Estructuras de Datos Avanzadas Notas de clase 01 Definici' on 16 (Subcadena LMS) Se define una subcadena LMS como una subcadenaS[i, j] tal que se cumpla alguna de las siguientes condiciones: i < jy tanto Suf(S, i) como Suf(S, j) son sufijos LMS y no existe un '  ndicei < k < j tal que Suf(S, k) sea un sufijo de tipo LMS. S[i, j] = Suf(S, n-1). En la...
\item[Pag. 20] CS3014 - Estructuras de Datos Avanzadas Notas de clase 01 Algoritmo 17:Compute-Buckets(n, S, end) 1Seacnt[0. . .(|Sigma| -1)] un arreglo de enteros inicializado con 0 ; 2parac in Shacer 3cnt[c]<-cnt[c] + 1 ; 4SeaP[0. . .(|Sigma| -1)] un arreglo de enteros ; 5sum<-0 ; 6parai<-0a|Sigma| -1hacer 7siend=F alseentonces 8P[i]<-sum; 9sum<-sum+cnt[i] ; 10en otro caso 11sum<-sum+cnt[i] ;...
\item[Pag. 21] CS3014 - Estructuras de Datos Avanzadas Notas de clase 01 Algoritmo 18:Induce-Order(n, S, T, X) 1SeaA[0. . .(n-1)] un arreglo de enteros inicializado en 0 ; 2R<-Compute-Buckets(n, S,True) ; 3parai<- |X| -1a0hacer 4c<-S[X[i]] ; 5A[R[c]]<-X[i] ; 6R[c]<-R[c]-1 ; 7L<-Compute-Buckets(n, S,F alse) ; 8parai<-0an-1hacer 9siA[i]>0andT[A[i]-1] =Lentonces 10c<-S[A[i]-1] ; 11A[...
\item[Pag. 22] CS3014 - Estructuras de Datos Avanzadas Notas de clase 01 esO(n). El siguiente teorema garantiza la correctitud del algoritmo SAIS. Teorema 23 (Nonget al., 2009) Es posible ordenar los sufijos deSllamando aInduce-Order(n, S, T,LMS suf f ix), donde LMS suf f ix esLMSordenado por el sufijo de cada posici' on correspondiente. Para calcularLMS suf f ix, primero...
\item[Pag. 23] CS3014 - Estructuras de Datos Avanzadas Notas de clase 01 [3] Udi Manber and Gene Myers. Suffix arrays: a new method for on-line string searches.siam Journal on Computing, 22(5):935-948, 1993. [4] Ge Nong, Sen Zhang, and Wai Hong Chan. Linear suffix array construction by almost pure induced-sorting. In2009 data compression conference, pages 193-202. IEEE, 20...
\end{description}
\subsection*{Estructuras\_de\_Datos\_Avanzadas\_\_\_Strings\_I.pdf (35 paginas)}
\begin{description}[leftmargin=!,labelwidth=4.2em,style=nextline]
\item[Pag. 1] CS3014 Estructuras de Datos Avanzadas CS3014 Estructuras de Datos Avanzadas Laboratorio - Semana 10 - Sesion 1 Victor Racso Galvan Oyola vgalvan@utec.edu.pe 27 de mayo de 2026
\item[Pag. 2] CS3014 Estructuras de Datos Avanzadas ?Que aprenderemos hoy? -RMQ y LCA -Pattern Matching
\item[Pag. 3] CS3014 Estructuras de Datos Avanzadas ?Que aprenderemos hoy? -RMQ y LCA -Pattern Matching
\item[Pag. 4] CS3014 Estructuras de Datos Avanzadas ?Que aprenderemos hoy? -RMQ y LCA -Pattern Matching
\item[Pag. 5] RMQ y LCA
\item[Pag. 6] CS3014 Estructuras de Datos Avanzadas RMQ Range Minimum Query (Codeforces Group - Static RMQ) Dada una secuenciaAdenenteros, calcular el valor m '  n l<=i<=r \{Ai \} para multiples consultas. Soluciones Implementaremos primero una version que tomeO(nlogn)de construccion yO(1)de consulta. Luego implementaremos una version que tomeO(n)de construccion y O(logn)de c...
\item[Pag. 7] CS3014 Estructuras de Datos Avanzadas RMQ Range Minimum Query (Codeforces Group - Static RMQ) Dada una secuenciaAdenenteros, calcular el valor m '  n l<=i<=r \{Ai \} para multiples consultas. Soluciones Implementaremos primero una version que tomeO(nlogn)de construccion yO(1)de consulta. Luego implementaremos una version que tomeO(n)de construccion y O(logn)de c...
\item[Pag. 8] CS3014 Estructuras de Datos Avanzadas RMQ Range Minimum Query (Codeforces Group - Static RMQ) Dada una secuenciaAdenenteros, calcular el valor m '  n l<=i<=r \{Ai \} para multiples consultas. Soluciones Implementaremos primero una version que tomeO(nlogn)de construccion yO(1)de consulta. Luego implementaremos una version que tomeO(n)de construccion y O(logn)de c...
\item[Pag. 9] CS3014 Estructuras de Datos Avanzadas RMQ Range Minimum Query (Codeforces Group - Static RMQ) Dada una secuenciaAdenenteros, calcular el valor m '  n l<=i<=r \{Ai \} para multiples consultas. Soluciones Implementaremos primero una version que tomeO(nlogn)de construccion yO(1)de consulta. Luego implementaremos una version que tomeO(n)de construccion y O(logn)de c...
\item[Pag. 10] CS3014 Estructuras de Datos Avanzadas RMQ Range Minimum Query (Codeforces Group - Static RMQ) Dada una secuenciaAdenenteros, calcular el valor m '  n l<=i<=r \{Ai \} para multiples consultas. Soluciones Implementaremos primero una version que tomeO(nlogn)de construccion yO(1)de consulta. Luego implementaremos una version que tomeO(n)de construccion y O(logn)de c...
\item[Pag. 11] CS3014 Estructuras de Datos Avanzadas RMQ Range Minimum Query (Codeforces Group - Static RMQ) Dada una secuenciaAdenenteros, calcular el valor m '  n l<=i<=r \{Ai \} para multiples consultas. Soluciones Implementaremos primero una version que tomeO(nlogn)de construccion yO(1)de consulta. Luego implementaremos una version que tomeO(n)de construccion y O(logn)de c...
\item[Pag. 12] CS3014 Estructuras de Datos Avanzadas LCA Lowest Common Ancestor Dado un arbol conNnodos y raiz en el nodo 0, responder aQconsultas (ui ,v i)con el LCA deu i yv i. Soluciones Implementaremos primero la solucion conO(NlogN)de construccion y O(logN)de consulta. Podemos implementar facilmente una solucion conO(N)de construccion y O(logN)de consulta. La implemen...
\item[Pag. 13] CS3014 Estructuras de Datos Avanzadas LCA Lowest Common Ancestor Dado un arbol conNnodos y raiz en el nodo 0, responder aQconsultas (ui ,v i)con el LCA deu i yv i. Soluciones Implementaremos primero la solucion conO(NlogN)de construccion y O(logN)de consulta. Podemos implementar facilmente una solucion conO(N)de construccion y O(logN)de consulta. La implemen...
\item[Pag. 14] CS3014 Estructuras de Datos Avanzadas LCA Lowest Common Ancestor Dado un arbol conNnodos y raiz en el nodo 0, responder aQconsultas (ui ,v i)con el LCA deu i yv i. Soluciones Implementaremos primero la solucion conO(NlogN)de construccion y O(logN)de consulta. Podemos implementar facilmente una solucion conO(N)de construccion y O(logN)de consulta. La implemen...
\item[Pag. 15] CS3014 Estructuras de Datos Avanzadas LCA Lowest Common Ancestor Dado un arbol conNnodos y raiz en el nodo 0, responder aQconsultas (ui ,v i)con el LCA deu i yv i. Soluciones Implementaremos primero la solucion conO(NlogN)de construccion y O(logN)de consulta. Podemos implementar facilmente una solucion conO(N)de construccion y O(logN)de consulta. La implemen...
\item[Pag. 16] CS3014 Estructuras de Datos Avanzadas LCA Lowest Common Ancestor Dado un arbol conNnodos y raiz en el nodo 0, responder aQconsultas (ui ,v i)con el LCA deu i yv i. Soluciones Implementaremos primero la solucion conO(NlogN)de construccion y O(logN)de consulta. Podemos implementar facilmente una solucion conO(N)de construccion y O(logN)de consulta. La implemen...
\item[Pag. 17] CS3014 Estructuras de Datos Avanzadas LCA Lowest Common Ancestor Dado un arbol conNnodos y raiz en el nodo 0, responder aQconsultas (ui ,v i)con el LCA deu i yv i. Soluciones Implementaremos primero la solucion conO(NlogN)de construccion y O(logN)de consulta. Podemos implementar facilmente una solucion conO(N)de construccion y O(logN)de consulta. La implemen...
\item[Pag. 18] Pattern Matching
\item[Pag. 19] CS3014 Estructuras de Datos Avanzadas Pattern Matching Definicion Es un problema en el que se tiene una cadenaT, llamada texto, y una cadena P, llamada patron, y se desea saber siPocurre enTcomo subcadena. Su variacion llamadaMultiple pattern matchingimplica tener multiples pa- tronesP 1, . . . ,Pm. Aplicaciones Busquedas exactas de texto, motores de busqued...
\item[Pag. 20] CS3014 Estructuras de Datos Avanzadas Pattern Matching Definicion Es un problema en el que se tiene una cadenaT, llamada texto, y una cadena P, llamada patron, y se desea saber siPocurre enTcomo subcadena. Su variacion llamadaMultiple pattern matchingimplica tener multiples pa- tronesP 1, . . . ,Pm. Aplicaciones Busquedas exactas de texto, motores de busqued...
\item[Pag. 21] CS3014 Estructuras de Datos Avanzadas Pattern Matching Definicion Es un problema en el que se tiene una cadenaT, llamada texto, y una cadena P, llamada patron, y se desea saber siPocurre enTcomo subcadena. Su variacion llamadaMultiple pattern matchingimplica tener multiples pa- tronesP 1, . . . ,Pm. Aplicaciones Busquedas exactas de texto, motores de busqued...
\item[Pag. 22] CS3014 Estructuras de Datos Avanzadas Pattern Matching Definicion Es un problema en el que se tiene una cadenaT, llamada texto, y una cadena P, llamada patron, y se desea saber siPocurre enTcomo subcadena. Su variacion llamadaMultiple pattern matchingimplica tener multiples pa- tronesP 1, . . . ,Pm. Aplicaciones Busquedas exactas de texto, motores de busqued...
\item[Pag. 23] CS3014 Estructuras de Datos Avanzadas Pattern Matching Definicion Es un problema en el que se tiene una cadenaT, llamada texto, y una cadena P, llamada patron, y se desea saber siPocurre enTcomo subcadena. Su variacion llamadaMultiple pattern matchingimplica tener multiples pa- tronesP 1, . . . ,Pm. Aplicaciones Busquedas exactas de texto, motores de busqued...
\item[Pag. 24] CS3014 Estructuras de Datos Avanzadas Pattern Matching Definicion Es un problema en el que se tiene una cadenaT, llamada texto, y una cadena P, llamada patron, y se desea saber siPocurre enTcomo subcadena. Su variacion llamadaMultiple pattern matchingimplica tener multiples pa- tronesP 1, . . . ,Pm. Aplicaciones Busquedas exactas de texto, motores de busqued...
\item[Pag. 25] CS3014 Estructuras de Datos Avanzadas Escenarios Pconocido,Tpor conocer Consideraremos un escenario en el que se nos daPinicialmente, procesaremos Py luego se nos daTpara hacer el emparejamiento. Algoritmos: Knuth-Morris-Pratt, Z, Rabin-Karp, Aho-Corasick (multiple). Tconocido,Ppor conocer Consideraremos un escenario en el que se nos daTinicialmente, procesa...
\item[Pag. 26] CS3014 Estructuras de Datos Avanzadas Escenarios Pconocido,Tpor conocer Consideraremos un escenario en el que se nos daPinicialmente, procesaremos Py luego se nos daTpara hacer el emparejamiento. Algoritmos: Knuth-Morris-Pratt, Z, Rabin-Karp, Aho-Corasick (multiple). Tconocido,Ppor conocer Consideraremos un escenario en el que se nos daTinicialmente, procesa...
\item[Pag. 27] CS3014 Estructuras de Datos Avanzadas Escenarios Pconocido,Tpor conocer Consideraremos un escenario en el que se nos daPinicialmente, procesaremos Py luego se nos daTpara hacer el emparejamiento. Algoritmos: Knuth-Morris-Pratt, Z, Rabin-Karp, Aho-Corasick (multiple). Tconocido,Ppor conocer Consideraremos un escenario en el que se nos daTinicialmente, procesa...
\item[Pag. 28] CS3014 Estructuras de Datos Avanzadas Escenarios Pconocido,Tpor conocer Consideraremos un escenario en el que se nos daPinicialmente, procesaremos Py luego se nos daTpara hacer el emparejamiento. Algoritmos: Knuth-Morris-Pratt, Z, Rabin-Karp, Aho-Corasick (multiple). Tconocido,Ppor conocer Consideraremos un escenario en el que se nos daTinicialmente, procesa...
\item[Pag. 29] CS3014 Estructuras de Datos Avanzadas Escenarios Pconocido,Tpor conocer Consideraremos un escenario en el que se nos daPinicialmente, procesaremos Py luego se nos daTpara hacer el emparejamiento. Algoritmos: Knuth-Morris-Pratt, Z, Rabin-Karp, Aho-Corasick (multiple). Tconocido,Ppor conocer Consideraremos un escenario en el que se nos daTinicialmente, procesa...
\item[Pag. 30] CS3014 Estructuras de Datos Avanzadas Escenarios Pconocido,Tpor conocer Consideraremos un escenario en el que se nos daPinicialmente, procesaremos Py luego se nos daTpara hacer el emparejamiento. Algoritmos: Knuth-Morris-Pratt, Z, Rabin-Karp, Aho-Corasick (multiple). Tconocido,Ppor conocer Consideraremos un escenario en el que se nos daTinicialmente, procesa...
\item[Pag. 31] CS3014 Estructuras de Datos Avanzadas Resumen de la sesion -Implementamos RMQ -Implementamos LCA -Aprendimos de que trata el problema de Pattern Matching
\item[Pag. 32] CS3014 Estructuras de Datos Avanzadas Resumen de la sesion -Implementamos RMQ -Implementamos LCA -Aprendimos de que trata el problema de Pattern Matching
\item[Pag. 33] CS3014 Estructuras de Datos Avanzadas Resumen de la sesion -Implementamos RMQ -Implementamos LCA -Aprendimos de que trata el problema de Pattern Matching
\item[Pag. 34] CS3014 Estructuras de Datos Avanzadas Resumen de la sesion -Implementamos RMQ -Implementamos LCA -Aprendimos de que trata el problema de Pattern Matching
\item[Pag. 35] CS3014 Estructuras de Datos Avanzadas Gracias
\end{description}
\subsection*{Estructuras\_de\_Datos\_Avanzadas\_\_\_Strings\_II.pdf (81 paginas)}
\begin{description}[leftmargin=!,labelwidth=4.2em,style=nextline]
\item[Pag. 1] CS3014 Estructuras de Datos Avanzadas CS3014 Estructuras de Datos Avanzadas Laboratorio - Semana 10 - Sesion 2 Victor Racso Galvan Oyola vgalvan@utec.edu.pe 29 de mayo de 2026
\item[Pag. 2] CS3014 Estructuras de Datos Avanzadas ?Que aprenderemos hoy? -Strings -Suffix array
\item[Pag. 3] CS3014 Estructuras de Datos Avanzadas ?Que aprenderemos hoy? -Strings -Suffix array
\item[Pag. 4] CS3014 Estructuras de Datos Avanzadas ?Que aprenderemos hoy? -Strings -Suffix array
\item[Pag. 5] Strings
\item[Pag. 6] CS3014 Estructuras de Datos Avanzadas Notacion Notacion basica de strings Una cadena es una secuencia de simbolos llamadoscaracteres, los cuales pertenecen a un conjunto llamadoalfabetoSigma. Eli-esimo caracter de una cadenaSse denota porS[i]y la concatenacion de los caracteres de las posicionesx in [i,j]se denotar porS[i,j]. Convenciones -Se considera comparacionl...
\item[Pag. 7] CS3014 Estructuras de Datos Avanzadas Notacion Notacion basica de strings Una cadena es una secuencia de simbolos llamadoscaracteres, los cuales pertenecen a un conjunto llamadoalfabetoSigma. Eli-esimo caracter de una cadenaSse denota porS[i]y la concatenacion de los caracteres de las posicionesx in [i,j]se denotar porS[i,j]. Convenciones -Se considera comparacionl...
\item[Pag. 8] CS3014 Estructuras de Datos Avanzadas Notacion Notacion basica de strings Una cadena es una secuencia de simbolos llamadoscaracteres, los cuales pertenecen a un conjunto llamadoalfabetoSigma. Eli-esimo caracter de una cadenaSse denota porS[i]y la concatenacion de los caracteres de las posicionesx in [i,j]se denotar porS[i,j]. Convenciones -Se considera comparacionl...
\item[Pag. 9] CS3014 Estructuras de Datos Avanzadas Notacion Notacion basica de strings Una cadena es una secuencia de simbolos llamadoscaracteres, los cuales pertenecen a un conjunto llamadoalfabetoSigma. Eli-esimo caracter de una cadenaSse denota porS[i]y la concatenacion de los caracteres de las posicionesx in [i,j]se denotar porS[i,j]. Convenciones -Se considera comparacionl...
\item[Pag. 10] CS3014 Estructuras de Datos Avanzadas Notacion Notacion basica de strings Una cadena es una secuencia de simbolos llamadoscaracteres, los cuales pertenecen a un conjunto llamadoalfabetoSigma. Eli-esimo caracter de una cadenaSse denota porS[i]y la concatenacion de los caracteres de las posicionesx in [i,j]se denotar porS[i,j]. Convenciones -Se considera comparacionl...
\item[Pag. 11] CS3014 Estructuras de Datos Avanzadas Notacion Notacion basica de strings Una cadena es una secuencia de simbolos llamadoscaracteres, los cuales pertenecen a un conjunto llamadoalfabetoSigma. Eli-esimo caracter de una cadenaSse denota porS[i]y la concatenacion de los caracteres de las posicionesx in [i,j]se denotar porS[i,j]. Convenciones -Se considera comparacionl...
\item[Pag. 12] Suffix array
\item[Pag. 13] CS3014 Estructuras de Datos Avanzadas Suffix array Definicion El suffix array de una cadenaScon longitudnes una permutacionAde sus posiciones de manera que siiva antes quejenA, entoncesS[i,n-1]< S[j,n-1]. Denotaremos porSuf(S,i)al sufijo deSque comienza en la posicioni, es decir,S[i,n-1]. Setup basico En primer lugar, para cualquier cadenaSque tengamos como...
\item[Pag. 14] CS3014 Estructuras de Datos Avanzadas Suffix array Definicion El suffix array de una cadenaScon longitudnes una permutacionAde sus posiciones de manera que siiva antes quejenA, entoncesS[i,n-1]< S[j,n-1]. Denotaremos porSuf(S,i)al sufijo deSque comienza en la posicioni, es decir,S[i,n-1]. Setup basico En primer lugar, para cualquier cadenaSque tengamos como...
\item[Pag. 15] CS3014 Estructuras de Datos Avanzadas Suffix array Definicion El suffix array de una cadenaScon longitudnes una permutacionAde sus posiciones de manera que siiva antes quejenA, entoncesS[i,n-1]< S[j,n-1]. Denotaremos porSuf(S,i)al sufijo deSque comienza en la posicioni, es decir,S[i,n-1]. Setup basico En primer lugar, para cualquier cadenaSque tengamos como...
\item[Pag. 16] CS3014 Estructuras de Datos Avanzadas Suffix array Definicion El suffix array de una cadenaScon longitudnes una permutacionAde sus posiciones de manera que siiva antes quejenA, entoncesS[i,n-1]< S[j,n-1]. Denotaremos porSuf(S,i)al sufijo deSque comienza en la posicioni, es decir,S[i,n-1]. Setup basico En primer lugar, para cualquier cadenaSque tengamos como...
\item[Pag. 17] CS3014 Estructuras de Datos Avanzadas Suffix array Definicion El suffix array de una cadenaScon longitudnes una permutacionAde sus posiciones de manera que siiva antes quejenA, entoncesS[i,n-1]< S[j,n-1]. Denotaremos porSuf(S,i)al sufijo deSque comienza en la posicioni, es decir,S[i,n-1]. Setup basico En primer lugar, para cualquier cadenaSque tengamos como...
\item[Pag. 18] CS3014 Estructuras de Datos Avanzadas Construccion del suffix array La siguiente tabla muestra algunos algoritmos de construccion de suffix array (SACA) importantes historicamente. Autor(es) Algoritmo Complejidad Idea Manber \& Myers Sin nombre O(nlogn) Doubling Itoh \& Tanaka Sin nombre Superlineal Ordenamiento inducido Ko \& Aluru SAKA O(n) Ordenamiento induc...
\item[Pag. 19] CS3014 Estructuras de Datos Avanzadas Construccion del suffix array La siguiente tabla muestra algunos algoritmos de construccion de suffix array (SACA) importantes historicamente. Autor(es) Algoritmo Complejidad Idea Manber \& Myers Sin nombre O(nlogn) Doubling Itoh \& Tanaka Sin nombre Superlineal Ordenamiento inducido Ko \& Aluru SAKA O(n) Ordenamiento induc...
\item[Pag. 20] CS3014 Estructuras de Datos Avanzadas DC3 (Karkkainen et. al.) Idea Particionamos las posiciones segun su residuo respecto a 3. Denotamos los tres conjuntos S0 =\{0<=i<=n-1:i=3k,para algunk in Z\} S1 =\{0<=i<=n-1:i=3k+1,para algunk in Z\} S2 =\{0<=i<=n-1:i=3k+2,para algunk in Z\} Teorema (Karkkainen et. al.) Si se sabe el orden de los sufijos en los elementos deS1 yS 2, entonces...
\item[Pag. 21] CS3014 Estructuras de Datos Avanzadas DC3 (Karkkainen et. al.) Idea Particionamos las posiciones segun su residuo respecto a 3. Denotamos los tres conjuntos S0 =\{0<=i<=n-1:i=3k,para algunk in Z\} S1 =\{0<=i<=n-1:i=3k+1,para algunk in Z\} S2 =\{0<=i<=n-1:i=3k+2,para algunk in Z\} Teorema (Karkkainen et. al.) Si se sabe el orden de los sufijos en los elementos deS1 yS 2, entonces...
\item[Pag. 22] CS3014 Estructuras de Datos Avanzadas DC3 (Karkkainen et. al.) Idea Particionamos las posiciones segun su residuo respecto a 3. Denotamos los tres conjuntos S0 =\{0<=i<=n-1:i=3k,para algunk in Z\} S1 =\{0<=i<=n-1:i=3k+1,para algunk in Z\} S2 =\{0<=i<=n-1:i=3k+2,para algunk in Z\} Teorema (Karkkainen et. al.) Si se sabe el orden de los sufijos en los elementos deS1 yS 2, entonces...
\item[Pag. 23] CS3014 Estructuras de Datos Avanzadas DC3 (Karkkainen et. al.) Idea Particionamos las posiciones segun su residuo respecto a 3. Denotamos los tres conjuntos S0 =\{0<=i<=n-1:i=3k,para algunk in Z\} S1 =\{0<=i<=n-1:i=3k+1,para algunk in Z\} S2 =\{0<=i<=n-1:i=3k+2,para algunk in Z\} Teorema (Karkkainen et. al.) Si se sabe el orden de los sufijos en los elementos deS1 yS 2, entonces...
\item[Pag. 24] CS3014 Estructuras de Datos Avanzadas DC3 (Karkkainen et. al.) Idea Particionamos las posiciones segun su residuo respecto a 3. Denotamos los tres conjuntos S0 =\{0<=i<=n-1:i=3k,para algunk in Z\} S1 =\{0<=i<=n-1:i=3k+1,para algunk in Z\} S2 =\{0<=i<=n-1:i=3k+2,para algunk in Z\} Teorema (Karkkainen et. al.) Si se sabe el orden de los sufijos en los elementos deS1 yS 2, entonces...
\item[Pag. 25] CS3014 Estructuras de Datos Avanzadas Algoritmo Particion ParaS 1 yS 2, vamos a tomar los bloques de 3 elementos en adelante desde cada posicion. Ya que podrian faltar elementos, se hace padding con sentinelas \$. Por ejemplo, para "banana\$", se tendrian los siguientes bloques: S1 ->(ana)(na\$) S2 ->(nan)(a\$\$) Notemos que es posible mapear cada uno de estos bloq...
\item[Pag. 26] CS3014 Estructuras de Datos Avanzadas Algoritmo Particion ParaS 1 yS 2, vamos a tomar los bloques de 3 elementos en adelante desde cada posicion. Ya que podrian faltar elementos, se hace padding con sentinelas \$. Por ejemplo, para "banana\$", se tendrian los siguientes bloques: S1 ->(ana)(na\$) S2 ->(nan)(a\$\$) Notemos que es posible mapear cada uno de estos bloq...
\item[Pag. 27] CS3014 Estructuras de Datos Avanzadas Algoritmo Particion ParaS 1 yS 2, vamos a tomar los bloques de 3 elementos en adelante desde cada posicion. Ya que podrian faltar elementos, se hace padding con sentinelas \$. Por ejemplo, para "banana\$", se tendrian los siguientes bloques: S1 ->(ana)(na\$) S2 ->(nan)(a\$\$) Notemos que es posible mapear cada uno de estos bloq...
\item[Pag. 28] CS3014 Estructuras de Datos Avanzadas Algoritmo Particion ParaS 1 yS 2, vamos a tomar los bloques de 3 elementos en adelante desde cada posicion. Ya que podrian faltar elementos, se hace padding con sentinelas \$. Por ejemplo, para "banana\$", se tendrian los siguientes bloques: S1 ->(ana)(na\$) S2 ->(nan)(a\$\$) Notemos que es posible mapear cada uno de estos bloq...
\item[Pag. 29] CS3014 Estructuras de Datos Avanzadas Algoritmo ?Y ahora que hacemos? Podemosreemplazarcadabloqueporsurespectivorankyseterminancreando dos cadenas nuevas. S1 ->R 1 =12 S2 ->R 2 =30 Si concatenamos estas dos cadenas nuevas, tendremos una cadenaR=R1R2 de longitud~ 2n 3 . Aplicando recursion Podemos usar recursion para obtener el suffix array de la cadenaR, de es...
\item[Pag. 30] CS3014 Estructuras de Datos Avanzadas Algoritmo ?Y ahora que hacemos? Podemosreemplazarcadabloqueporsurespectivorankyseterminancreando dos cadenas nuevas. S1 ->R 1 =12 S2 ->R 2 =30 Si concatenamos estas dos cadenas nuevas, tendremos una cadenaR=R1R2 de longitud~ 2n 3 . Aplicando recursion Podemos usar recursion para obtener el suffix array de la cadenaR, de es...
\item[Pag. 31] CS3014 Estructuras de Datos Avanzadas Algoritmo ?Y ahora que hacemos? Podemosreemplazarcadabloqueporsurespectivorankyseterminancreando dos cadenas nuevas. S1 ->R 1 =12 S2 ->R 2 =30 Si concatenamos estas dos cadenas nuevas, tendremos una cadenaR=R1R2 de longitud~ 2n 3 . Aplicando recursion Podemos usar recursion para obtener el suffix array de la cadenaR, de es...
\item[Pag. 32] CS3014 Estructuras de Datos Avanzadas Algoritmo ?Y ahora que hacemos? Podemosreemplazarcadabloqueporsurespectivorankyseterminancreando dos cadenas nuevas. S1 ->R 1 =12 S2 ->R 2 =30 Si concatenamos estas dos cadenas nuevas, tendremos una cadenaR=R1R2 de longitud~ 2n 3 . Aplicando recursion Podemos usar recursion para obtener el suffix array de la cadenaR, de es...
\item[Pag. 33] CS3014 Estructuras de Datos Avanzadas Algoritmo ?Y ahora que hacemos? Podemosreemplazarcadabloqueporsurespectivorankyseterminancreando dos cadenas nuevas. S1 ->R 1 =12 S2 ->R 2 =30 Si concatenamos estas dos cadenas nuevas, tendremos una cadenaR=R1R2 de longitud~ 2n 3 . Aplicando recursion Podemos usar recursion para obtener el suffix array de la cadenaR, de es...
\item[Pag. 34] CS3014 Estructuras de Datos Avanzadas Algoritmo ?Y ahora que hacemos? Podemosreemplazarcadabloqueporsurespectivorankyseterminancreando dos cadenas nuevas. S1 ->R 1 =12 S2 ->R 2 =30 Si concatenamos estas dos cadenas nuevas, tendremos una cadenaR=R1R2 de longitud~ 2n 3 . Aplicando recursion Podemos usar recursion para obtener el suffix array de la cadenaR, de es...
\item[Pag. 35] CS3014 Estructuras de Datos Avanzadas Algoritmo ?Y como ordeno todo? Vamos a tener escenarios en los cuales todo se va a reducir a comparar una cantidadO(1)de caracteres y una comparacion de ranks de los sufijos deS1 yS 2. -x in  S 0 vsy in  S 1: Podemos tomar el primer caracter de cada sufijo, de manera que la comparacion se reduce a comparar los pares (S[i],rank...
\item[Pag. 36] CS3014 Estructuras de Datos Avanzadas Algoritmo ?Y como ordeno todo? Vamos a tener escenarios en los cuales todo se va a reducir a comparar una cantidadO(1)de caracteres y una comparacion de ranks de los sufijos deS1 yS 2. -x in  S 0 vsy in  S 1: Podemos tomar el primer caracter de cada sufijo, de manera que la comparacion se reduce a comparar los pares (S[i],rank...
\item[Pag. 37] CS3014 Estructuras de Datos Avanzadas Algoritmo ?Y como ordeno todo? Vamos a tener escenarios en los cuales todo se va a reducir a comparar una cantidadO(1)de caracteres y una comparacion de ranks de los sufijos deS1 yS 2. -x in  S 0 vsy in  S 1: Podemos tomar el primer caracter de cada sufijo, de manera que la comparacion se reduce a comparar los pares (S[i],rank...
\item[Pag. 38] CS3014 Estructuras de Datos Avanzadas Complejidad Analisis Ya que es posible hacer la comparacion de pares y tuplas enO(n)usando radix sort, obtenemos una recurrencia de la siguiente forma: T(n) =T  2n 3   +O(n) Es sencillo notar queT(n) =O(n). Generalizacion Se planteo en investigaciones posteriores la posibilidad de tomar modulok en vez de modulo 3. El mod...
\item[Pag. 39] CS3014 Estructuras de Datos Avanzadas Complejidad Analisis Ya que es posible hacer la comparacion de pares y tuplas enO(n)usando radix sort, obtenemos una recurrencia de la siguiente forma: T(n) =T  2n 3   +O(n) Es sencillo notar queT(n) =O(n). Generalizacion Se planteo en investigaciones posteriores la posibilidad de tomar modulok en vez de modulo 3. El mod...
\item[Pag. 40] CS3014 Estructuras de Datos Avanzadas Complejidad Analisis Ya que es posible hacer la comparacion de pares y tuplas enO(n)usando radix sort, obtenemos una recurrencia de la siguiente forma: T(n) =T  2n 3   +O(n) Es sencillo notar queT(n) =O(n). Generalizacion Se planteo en investigaciones posteriores la posibilidad de tomar modulok en vez de modulo 3. El mod...
\item[Pag. 41] CS3014 Estructuras de Datos Avanzadas Complejidad Analisis Ya que es posible hacer la comparacion de pares y tuplas enO(n)usando radix sort, obtenemos una recurrencia de la siguiente forma: T(n) =T  2n 3   +O(n) Es sencillo notar queT(n) =O(n). Generalizacion Se planteo en investigaciones posteriores la posibilidad de tomar modulok en vez de modulo 3. El mod...
\item[Pag. 42] CS3014 Estructuras de Datos Avanzadas Complejidad Analisis Ya que es posible hacer la comparacion de pares y tuplas enO(n)usando radix sort, obtenemos una recurrencia de la siguiente forma: T(n) =T  2n 3   +O(n) Es sencillo notar queT(n) =O(n). Generalizacion Se planteo en investigaciones posteriores la posibilidad de tomar modulok en vez de modulo 3. El mod...
\item[Pag. 43] CS3014 Estructuras de Datos Avanzadas Complejidad Analisis Ya que es posible hacer la comparacion de pares y tuplas enO(n)usando radix sort, obtenemos una recurrencia de la siguiente forma: T(n) =T  2n 3   +O(n) Es sencillo notar queT(n) =O(n). Generalizacion Se planteo en investigaciones posteriores la posibilidad de tomar modulok en vez de modulo 3. El mod...
\item[Pag. 44] CS3014 Estructuras de Datos Avanzadas Complejidad Analisis Ya que es posible hacer la comparacion de pares y tuplas enO(n)usando radix sort, obtenemos una recurrencia de la siguiente forma: T(n) =T  2n 3   +O(n) Es sencillo notar queT(n) =O(n). Generalizacion Se planteo en investigaciones posteriores la posibilidad de tomar modulok en vez de modulo 3. El mod...
\item[Pag. 45] CS3014 Estructuras de Datos Avanzadas SAIS (Nong et. al.) Idea Clasificar los sufijos deSen tipoSyL. Un sufijoSuf(S,i)es de tipo -S, siSuf(S,i)<Suf(S,i+1). -L, siSuf(S,i)>Suf(S,i+1). Un sufijoSuf(S,i)es de tipo LMS si el sufijoSuf(S,i-1)es de tipoLy el sufijoSuf(S,i)es de tipoS. Notemos que hay a lo muchon-1 2 sufijos de tipo LMS. Teorema (Nong et. al.) Si s...
\item[Pag. 46] CS3014 Estructuras de Datos Avanzadas SAIS (Nong et. al.) Idea Clasificar los sufijos deSen tipoSyL. Un sufijoSuf(S,i)es de tipo -S, siSuf(S,i)<Suf(S,i+1). -L, siSuf(S,i)>Suf(S,i+1). Un sufijoSuf(S,i)es de tipo LMS si el sufijoSuf(S,i-1)es de tipoLy el sufijoSuf(S,i)es de tipoS. Notemos que hay a lo muchon-1 2 sufijos de tipo LMS. Teorema (Nong et. al.) Si s...
\item[Pag. 47] CS3014 Estructuras de Datos Avanzadas SAIS (Nong et. al.) Idea Clasificar los sufijos deSen tipoSyL. Un sufijoSuf(S,i)es de tipo -S, siSuf(S,i)<Suf(S,i+1). -L, siSuf(S,i)>Suf(S,i+1). Un sufijoSuf(S,i)es de tipo LMS si el sufijoSuf(S,i-1)es de tipoLy el sufijoSuf(S,i)es de tipoS. Notemos que hay a lo muchon-1 2 sufijos de tipo LMS. Teorema (Nong et. al.) Si s...
\item[Pag. 48] CS3014 Estructuras de Datos Avanzadas SAIS (Nong et. al.) Idea Clasificar los sufijos deSen tipoSyL. Un sufijoSuf(S,i)es de tipo -S, siSuf(S,i)<Suf(S,i+1). -L, siSuf(S,i)>Suf(S,i+1). Un sufijoSuf(S,i)es de tipo LMS si el sufijoSuf(S,i-1)es de tipoLy el sufijoSuf(S,i)es de tipoS. Notemos que hay a lo muchon-1 2 sufijos de tipo LMS. Teorema (Nong et. al.) Si s...
\item[Pag. 49] CS3014 Estructuras de Datos Avanzadas SAIS (Nong et. al.) Idea Clasificar los sufijos deSen tipoSyL. Un sufijoSuf(S,i)es de tipo -S, siSuf(S,i)<Suf(S,i+1). -L, siSuf(S,i)>Suf(S,i+1). Un sufijoSuf(S,i)es de tipo LMS si el sufijoSuf(S,i-1)es de tipoLy el sufijoSuf(S,i)es de tipoS. Notemos que hay a lo muchon-1 2 sufijos de tipo LMS. Teorema (Nong et. al.) Si s...
\item[Pag. 50] CS3014 Estructuras de Datos Avanzadas SAIS (Nong et. al.) Idea Clasificar los sufijos deSen tipoSyL. Un sufijoSuf(S,i)es de tipo -S, siSuf(S,i)<Suf(S,i+1). -L, siSuf(S,i)>Suf(S,i+1). Un sufijoSuf(S,i)es de tipo LMS si el sufijoSuf(S,i-1)es de tipoLy el sufijoSuf(S,i)es de tipoS. Notemos que hay a lo muchon-1 2 sufijos de tipo LMS. Teorema (Nong et. al.) Si s...
\item[Pag. 51] CS3014 Estructuras de Datos Avanzadas Algoritmo Particion Se realiza la particion deStomando en cuenta las posiciones de los sufijos LMS y se ordenan segun(caracter,tipo)para calcular el rank. Es posible realizar esto enO(n). Reduccion Se obtiene una cadena reducidaRen base a la particion por LMS, se calcula el suffix array deRy se puede inducir el orden de...
\item[Pag. 52] CS3014 Estructuras de Datos Avanzadas Algoritmo Particion Se realiza la particion deStomando en cuenta las posiciones de los sufijos LMS y se ordenan segun(caracter,tipo)para calcular el rank. Es posible realizar esto enO(n). Reduccion Se obtiene una cadena reducidaRen base a la particion por LMS, se calcula el suffix array deRy se puede inducir el orden de...
\item[Pag. 53] CS3014 Estructuras de Datos Avanzadas Algoritmo Particion Se realiza la particion deStomando en cuenta las posiciones de los sufijos LMS y se ordenan segun(caracter,tipo)para calcular el rank. Es posible realizar esto enO(n). Reduccion Se obtiene una cadena reducidaRen base a la particion por LMS, se calcula el suffix array deRy se puede inducir el orden de...
\item[Pag. 54] CS3014 Estructuras de Datos Avanzadas Algoritmo Particion Se realiza la particion deStomando en cuenta las posiciones de los sufijos LMS y se ordenan segun(caracter,tipo)para calcular el rank. Es posible realizar esto enO(n). Reduccion Se obtiene una cadena reducidaRen base a la particion por LMS, se calcula el suffix array deRy se puede inducir el orden de...
\item[Pag. 55] CS3014 Estructuras de Datos Avanzadas Algoritmo Particion Se realiza la particion deStomando en cuenta las posiciones de los sufijos LMS y se ordenan segun(caracter,tipo)para calcular el rank. Es posible realizar esto enO(n). Reduccion Se obtiene una cadena reducidaRen base a la particion por LMS, se calcula el suffix array deRy se puede inducir el orden de...
\item[Pag. 56] CS3014 Estructuras de Datos Avanzadas Algoritmo Particion Se realiza la particion deStomando en cuenta las posiciones de los sufijos LMS y se ordenan segun(caracter,tipo)para calcular el rank. Es posible realizar esto enO(n). Reduccion Se obtiene una cadena reducidaRen base a la particion por LMS, se calcula el suffix array deRy se puede inducir el orden de...
\item[Pag. 57] CS3014 Estructuras de Datos Avanzadas Calculando el LCP LCP array Es un arreglo en el cual se almacenan los Longest Common Prefix (LCP) entre cada par de posiciones adyacentes del suffix array Ventajas -Ya queaesta ordenado, nos permite calcular el LCP entre cualquier par de posiciones en un tiempo prudente si nos apoyamos en alguna estructura de datos. -Ya...
\item[Pag. 58] CS3014 Estructuras de Datos Avanzadas Calculando el LCP LCP array Es un arreglo en el cual se almacenan los Longest Common Prefix (LCP) entre cada par de posiciones adyacentes del suffix array Ventajas -Ya queaesta ordenado, nos permite calcular el LCP entre cualquier par de posiciones en un tiempo prudente si nos apoyamos en alguna estructura de datos. -Ya...
\item[Pag. 59] CS3014 Estructuras de Datos Avanzadas Calculando el LCP LCP array Es un arreglo en el cual se almacenan los Longest Common Prefix (LCP) entre cada par de posiciones adyacentes del suffix array Ventajas -Ya queaesta ordenado, nos permite calcular el LCP entre cualquier par de posiciones en un tiempo prudente si nos apoyamos en alguna estructura de datos. -Ya...
\item[Pag. 60] CS3014 Estructuras de Datos Avanzadas Calculando el LCP LCP array Es un arreglo en el cual se almacenan los Longest Common Prefix (LCP) entre cada par de posiciones adyacentes del suffix array Ventajas -Ya queaesta ordenado, nos permite calcular el LCP entre cualquier par de posiciones en un tiempo prudente si nos apoyamos en alguna estructura de datos. -Ya...
\item[Pag. 61] CS3014 Estructuras de Datos Avanzadas Calculando el LCP LCP array Es un arreglo en el cual se almacenan los Longest Common Prefix (LCP) entre cada par de posiciones adyacentes del suffix array Ventajas -Ya queaesta ordenado, nos permite calcular el LCP entre cualquier par de posiciones en un tiempo prudente si nos apoyamos en alguna estructura de datos. -Ya...
\item[Pag. 62] CS3014 Estructuras de Datos Avanzadas Calculando el LCP ?Y como se calcula? Una implementacion ingenua podria tomar hastaO(n2)de complejidad, asi que usaremos el algoritmo de Kasai para calcularlo en un mejor tiempo. Algoritmo de Kasai Su principal observacion es que si tenemos dos sufijos que son adyacentes en a, seanxyysus posiciones iniciales, con LCP igu...
\item[Pag. 63] CS3014 Estructuras de Datos Avanzadas Calculando el LCP ?Y como se calcula? Una implementacion ingenua podria tomar hastaO(n2)de complejidad, asi que usaremos el algoritmo de Kasai para calcularlo en un mejor tiempo. Algoritmo de Kasai Su principal observacion es que si tenemos dos sufijos que son adyacentes en a, seanxyysus posiciones iniciales, con LCP igu...
\item[Pag. 64] CS3014 Estructuras de Datos Avanzadas Calculando el LCP ?Y como se calcula? Una implementacion ingenua podria tomar hastaO(n2)de complejidad, asi que usaremos el algoritmo de Kasai para calcularlo en un mejor tiempo. Algoritmo de Kasai Su principal observacion es que si tenemos dos sufijos que son adyacentes en a, seanxyysus posiciones iniciales, con LCP igu...
\item[Pag. 65] CS3014 Estructuras de Datos Avanzadas Calculando el LCP ?Y como se calcula? Una implementacion ingenua podria tomar hastaO(n2)de complejidad, asi que usaremos el algoritmo de Kasai para calcularlo en un mejor tiempo. Algoritmo de Kasai Su principal observacion es que si tenemos dos sufijos que son adyacentes en a, seanxyysus posiciones iniciales, con LCP igu...
\item[Pag. 66] CS3014 Estructuras de Datos Avanzadas Calculando el LCP ?Y como se calcula? Una implementacion ingenua podria tomar hastaO(n2)de complejidad, asi que usaremos el algoritmo de Kasai para calcularlo en un mejor tiempo. Algoritmo de Kasai Su principal observacion es que si tenemos dos sufijos que son adyacentes en a, seanxyysus posiciones iniciales, con LCP igu...
\item[Pag. 67] CS3014 Estructuras de Datos Avanzadas Calculando el LCP Pero podrian no ser adyacentes Eso es cierto, pero si su LCP es al menosk-1, eso implica que todas las posiciones entre ellos tambien tienen LCP de al menosk-1 con cualquiera de los dos. Complejidad Sikse reduce en 1 cada vez que comparamos un par de posiciones dea, este es reducido a lo muchon-1 veces...
\item[Pag. 68] CS3014 Estructuras de Datos Avanzadas Calculando el LCP Pero podrian no ser adyacentes Eso es cierto, pero si su LCP es al menosk-1, eso implica que todas las posiciones entre ellos tambien tienen LCP de al menosk-1 con cualquiera de los dos. Complejidad Sikse reduce en 1 cada vez que comparamos un par de posiciones dea, este es reducido a lo muchon-1 veces...
\item[Pag. 69] CS3014 Estructuras de Datos Avanzadas Calculando el LCP Pero podrian no ser adyacentes Eso es cierto, pero si su LCP es al menosk-1, eso implica que todas las posiciones entre ellos tambien tienen LCP de al menosk-1 con cualquiera de los dos. Complejidad Sikse reduce en 1 cada vez que comparamos un par de posiciones dea, este es reducido a lo muchon-1 veces...
\item[Pag. 70] CS3014 Estructuras de Datos Avanzadas Calculando el LCP Pero podrian no ser adyacentes Eso es cierto, pero si su LCP es al menosk-1, eso implica que todas las posiciones entre ellos tambien tienen LCP de al menosk-1 con cualquiera de los dos. Complejidad Sikse reduce en 1 cada vez que comparamos un par de posiciones dea, este es reducido a lo muchon-1 veces...
\item[Pag. 71] CS3014 Estructuras de Datos Avanzadas Calculando el LCP Suffix Index LCP \$ 6 - A\$ 5 0 ANA\$ 3 1 ANANA\$ 1 3 BANANA\$ 0 0 NA\$ 4 0 NANA\$ 2 2
\item[Pag. 72] CS3014 Estructuras de Datos Avanzadas Propiedades RMQ Se puede calcular el LCP entre un par de sufijos deSmediante una consulta de RMQ sobre el LCP array. (Spoiler) El suffix tree es equivalente al auxiliary/tree de todos los sufijos y el LCA representa al LCP entre cada par de hojas. Conclusiones Los suffix array son estructuras poderosas que permiten obten...
\item[Pag. 73] CS3014 Estructuras de Datos Avanzadas Propiedades RMQ Se puede calcular el LCP entre un par de sufijos deSmediante una consulta de RMQ sobre el LCP array. (Spoiler) El suffix tree es equivalente al auxiliary/tree de todos los sufijos y el LCA representa al LCP entre cada par de hojas. Conclusiones Los suffix array son estructuras poderosas que permiten obten...
\item[Pag. 74] CS3014 Estructuras de Datos Avanzadas Propiedades RMQ Se puede calcular el LCP entre un par de sufijos deSmediante una consulta de RMQ sobre el LCP array. (Spoiler) El suffix tree es equivalente al auxiliary/tree de todos los sufijos y el LCA representa al LCP entre cada par de hojas. Conclusiones Los suffix array son estructuras poderosas que permiten obten...
\item[Pag. 75] CS3014 Estructuras de Datos Avanzadas Propiedades RMQ Se puede calcular el LCP entre un par de sufijos deSmediante una consulta de RMQ sobre el LCP array. (Spoiler) El suffix tree es equivalente al auxiliary/tree de todos los sufijos y el LCA representa al LCP entre cada par de hojas. Conclusiones Los suffix array son estructuras poderosas que permiten obten...
\item[Pag. 76] CS3014 Estructuras de Datos Avanzadas Propiedades RMQ Se puede calcular el LCP entre un par de sufijos deSmediante una consulta de RMQ sobre el LCP array. (Spoiler) El suffix tree es equivalente al auxiliary/tree de todos los sufijos y el LCA representa al LCP entre cada par de hojas. Conclusiones Los suffix array son estructuras poderosas que permiten obten...
\item[Pag. 77] CS3014 Estructuras de Datos Avanzadas Resumen de la sesion -Aprendimos suffix array -Aprendimos el algoritmo DC3 y SAIS -Aprendimos LCP array
\item[Pag. 78] CS3014 Estructuras de Datos Avanzadas Resumen de la sesion -Aprendimos suffix array -Aprendimos el algoritmo DC3 y SAIS -Aprendimos LCP array
\item[Pag. 79] CS3014 Estructuras de Datos Avanzadas Resumen de la sesion -Aprendimos suffix array -Aprendimos el algoritmo DC3 y SAIS -Aprendimos LCP array
\item[Pag. 80] CS3014 Estructuras de Datos Avanzadas Resumen de la sesion -Aprendimos suffix array -Aprendimos el algoritmo DC3 y SAIS -Aprendimos LCP array
\item[Pag. 81] CS3014 Estructuras de Datos Avanzadas Gracias
\end{description}
\subsection*{eda\_slides\_08\_rmq.pdf (27 paginas)}
\begin{description}[leftmargin=!,labelwidth=4.2em,style=nextline]
\item[Pag. 1] Static T rees CS3014 - Estructura de Datos Avanzados Luciano A. Romero Calla lromeroc@utec.edu.pe 2026-1
\item[Pag. 2] Overview 1. Goals 2. RMQ 3. RMQ \& LCA 4. Solving RMQ 5. Level Ancestor 6. Related 7. Summary 2 / 20
\item[Pag. 3] Goals 1. Goals 3 / 20
\item[Pag. 4] Goals Goals Warming up! Problem I - Latin American Regional ACM ICPC 2017 https://maratona.sbc.org.br/hist/2017/resultados/contest.pdf Do it yourself! (homework)https://judge.beecrowd.com/en/problems/view/2703 4 / 20
\item[Pag. 5] Goals Goals Warming up! Problem I - Latin American Regional ACM ICPC 2017 https://maratona.sbc.org.br/hist/2017/resultados/contest.pdf Do it yourself! (homework)https://judge.beecrowd.com/en/problems/view/2703 4 / 20
\item[Pag. 6] Goals Goals - Solve the RMQ problem query inO(1) - Explore different structures and analyze their performance. - Be able to reduce the LCA problem to RMQ and solve it. - Solve the Level Ancestor (LA) problem. 5 / 20
\item[Pag. 7] Goals Goals - Solve the RMQ problem query inO(1) - Explore different structures and analyze their performance. - Be able to reduce the LCA problem to RMQ and solve it. - Solve the Level Ancestor (LA) problem. 5 / 20
\item[Pag. 8] Goals Goals - Solve the RMQ problem query inO(1) - Explore different structures and analyze their performance. - Be able to reduce the LCA problem to RMQ and solve it. - Solve the Level Ancestor (LA) problem. 5 / 20
\item[Pag. 9] Goals Goals - Solve the RMQ problem query inO(1) - Explore different structures and analyze their performance. - Be able to reduce the LCA problem to RMQ and solve it. - Solve the Level Ancestor (LA) problem. 5 / 20
\item[Pag. 10] Goals Goals Problem Scenario Preprocessing Space Query Time Range Minimum Query (RMQ)O(n)O(n)O(1) Lowest Common Ancestor (LCA)O(n)O(n)O(1) Level Ancestor (LA)O(n)O(n)O(1) Goal: Achieve optimal linear preprocessing and constant time queries across all structures. 6 / 20
\item[Pag. 11] RMQ 2. RMQ 7 / 20
\item[Pag. 12] RMQ RMQ Problem Given an arrayAof n numbers (to preprocess). In a query, the goal is to find the minimum element in a range spanned byA[i]andA[j]: RMQ(i,j) = arg min\{A[i],A[i+1], ...,A[j]\}=k wherei<=k<=jandA[k]is minimized 8 / 20
\item[Pag. 13] RMQ \& LCA 3. RMQ \& LCA 3.1 Cartesian trees 3.2 LCA to 1RMQ: Euler Tour 9 / 20
\item[Pag. 14] RMQ \& LCA Cartesian trees Cartesian trees Rules: - Root = Minimum elementA[m] - Left / Right subtrees built recursively. - Preservesmin-heap property. Identity RMQA(i,j)=LCA T (i,j) A= [7,8,2,6,9,4] 2 7 8 4 6 9 10 / 20
\item[Pag. 15] RMQ \& LCA Cartesian trees Cartesian trees Linear time algorithm - Scan array left to right. - Walk up theright spineuntil findingA[v]<A[i]. - InsertA[i]as right child; old subtree becomes left child ofA[i]. - Amortized Cost:Nodes enter/leave spine <= 1 time ==>O(n). 2 4 7 5 11 / 20
\item[Pag. 16] RMQ \& LCA LCA to 1RMQ: Euler T our LCA to 1RMQ: Euler T our A B D C Euler Tour (E):A B D B A C A Depth Array (D): 0 1 2 1 0 1 0 The 1Property Consecutive depth steps change by exactly+1 or-1: |D[i]-D[i-1]|=1 LCA(u,v) = min value inDbetween indices ofuandv. 12 / 20
\item[Pag. 17] Solving RMQ 4. Solving RMQ 13 / 20
\item[Pag. 18] Solving RMQ Solving RMQ Block partition size:b= 1 2 log2 n. 1. Macro-Structure (Across Blocks) - Sparse Table over block minimums. - Space:O( n b log n b ) =O(n). 2. Micro-Structure (Inside Blocks) - Max unique shapes = 2b-1 =O(  n). - Precomputebxblookup tables. - Space:O(  n.b 2) =o(n). Total Performance:O(n)Space / O(1)Query Time 14 / 20
\item[Pag. 19] Level Ancestor 5. Level Ancestor 15 / 20
\item[Pag. 20] Level Ancestor Level Ancestor Strategy / Name Preprocessing Query Time Core T echnique Full Table LookupO(n 2)O(1)Complete DP lookup table Jump PointersO(nlogn)O(logn)Binary lifting (powers of 2) Longest Path DecompositionO(n)O(  n)Disjoint root-to-leaf path arrays Ladder DecompositionO(n)O(logn)Overlapping path extensions (doubled length) Jump + LadderO(nlo...
\item[Pag. 21] Related 6. Related 17 / 20
\item[Pag. 22] Related Related Dynamic Segment tree, Fenwick tree, ... 18 / 20
\item[Pag. 23] Summary 7. Summary 19 / 20
\item[Pag. 24] Summary Summary Problem Scenario Preprocessing Space Query Time Range Minimum Query (RMQ)O(n)O(n)O(1) Lowest Common Ancestor (LCA)O(n)O(n)O(1) Level Ancestor (LA)O(n)O(n)O(lgn),O(1) Goal: Achieve optimal linear preprocessing and constant time queries across all structures. 20 / 20
\item[Pag. 25] Thanks
\item[Pag. 26] References References 22 / 20
\item[Pag. 27] Acknowledgements Acknowledgements The course slides are based on the lectures from previous editions and on similar lectures elsewhere. List of credits: Erick Demaine, Keith Schwarz 23 / 20
\end{description}
